\section{Global comparison and the equality case}
\label{rig:section}

We retain all intermediate model types and the normalization of
Section~\ref{geom:section}. The preceding geodesic comparison
now gives the global volume bound. Its nonnegative defect will
also retain enough information to identify the equality metric.

For $\xi\in U$, write $\gamma_\xi$ for the geodesic with initial
velocity $\xi$ and $f_\xi$ for the determinant of
Proposition~\ref{lag:determinant}. Put
\[
 J(\xi,t)=f_\xi(0,t).
\]
This is the absolute determinant of the normal Jacobi tensor
with initial value zero and derivative $\Id$. It is continuous
in $(\xi,t)$, positive for $0<t<D$, and agrees there with the
usual polar volume density.

\begin{lemma}[Polar integration over a full period]
\label{rig:polar}
For every $p\in M$,
\[
 \int_{S_pM}\int_0^D J(\xi,t)\,dt\,d\sigma_p(\xi)=V,
 \qquad
 \int_{S_pM}\int_0^T J(\xi,t)\,dt\,d\sigma_p(\xi)=2V.
\]
\end{lemma}

\begin{proof}
The open injectivity ball maps diffeomorphically onto the
complement of the cut locus. The latter is a smooth submanifold
of dimension $d-k<n$, hence has zero Riemannian volume.
Polar coordinates therefore prove the first formula.
Periodicity of the Jacobi data and reversal of the geodesic give
\[
 J(\xi,T-t)=J(\xi,-t)=J(-\xi,t).
\]
The sign in the reversed initial derivative disappears under the
absolute determinant. The substitution $t\mapsto T-t$ in the
second half-period, followed by the area-preserving antipodal
map on $S_pM$, proves the second formula.
\end{proof}

\begin{proof}[Proof of Theorem~\ref{thm:comparison}]
The circle, sphere-type, and real-projective-type cases are
already covered in Section~\ref{geom:section}. In an intermediate
case use the normalization \eqref{geom:normalization}.
Theorem~\ref{weighted:geodesic} says that
\begin{equation}\label{rig:geodesic-defect}
 \delta(\xi):=\int_0^T\int_0^T
   \bigl(f_\xi(x,y)-w(|y-x|)\bigr)\,dx\,dy\ge0.
\end{equation}
Time translation along a geodesic gives
\[
 f_\xi(x,y)=J\bigl(\phi_x\xi,(y-x)\bmod T\bigr).
\]
The possible choice of endpoint for the residue modulo $T$ has
no effect on any integral. Fubini's theorem, translation on the
time circle, flow invariance of $\mu$, and
Lemma~\ref{rig:polar} yield
\begin{align*}
 \int_U\int_0^T\int_0^T f_\xi(x,y)\,dx\,dy\,d\mu(\xi)
 &=T\int_U\int_0^T J(\xi,t)\,dt\,d\mu(\xi)\\
 &=2TV^2.
\end{align*}
All integrands are continuous and bounded on the compact
domains in question. Since $w(T-t)=w(t)$, its periodic
extension gives likewise
\begin{align*}
 \int_U\int_0^T\int_0^T w(|y-x|)\,dx\,dy\,d\mu(\xi)
 &=\sigma_d V T\int_0^T w(t)\,dt\\
 &=2TVV_0.
\end{align*}
Integrating \eqref{rig:geodesic-defect} and dividing by $2TV>0$
proves $V\ge V_0$.

If $V=V_0$, the nonnegative function $\delta$ has integral zero.
It is continuous, because the underlying Jacobi tensors depend
smoothly on the initial velocity and taking the absolute
determinant preserves continuity. Full support of $\mu$ forces
$\delta(\xi)=0$ for every $\xi$. The equality assertion in
Theorem~\ref{weighted:geodesic} now gives
\begin{equation}\label{rig:radial-density}
 J(\xi,t)=w(t),\qquad \xi\in U,\quad 0<t<D.
\end{equation}

In Cartesian normal coordinates at any point, the volume
density is $J(\xi,t)/t^d$. By \eqref{rig:radial-density} it depends
only on $t$. At the origin it extends smoothly with value one,
as also follows directly from the normal-coordinate volume
form. Thus $M$ is harmonic. It is compact and simply connected
by Proposition~\ref{geom:structure}. The compact simply
connected case of the Lichnerowicz conjecture, proved by
\citet{Szabo1990}, identifies its metric with a compact rank-one
symmetric metric up to scale. Its type identifies the model,
and the fixed diameter identifies the scale.
For the radial-density definition and the concluding metric
classification, see also the author's original preprint
\citep[pp.~2--4 and Lemmas~4.9--4.10]{Szabo1988}.

Conversely, a model is isometric to itself and has volume $V_0$.
Finally both sides of the volume comparison scale by the same
factor under a constant metric rescaling. Undoing
\eqref{geom:normalization} proves the assertion at every diameter.
\end{proof}

\begin{proof}[Proof of Corollary~\ref{thm:volume-reduction}]
The upper bound and Theorem~\ref{thm:comparison} give equal
volumes, so its equality case applies.
\end{proof}

